\section{\texttt{run.py}: function-by-function implementation}
This module assembles the normal cipher traces and attack functions into terminal lessons. Compatibility wrappers delegate to piret rather than copying equations. The combined main returns success only when each selected full-recovery lesson succeeds.
\subsection{\texttt{block}}
\begin{lightcode}
def block(value):
\end{lightcode}
\noindent\textbf{Source:} complete \texttt{run.py} listing in Part 10, lines 10--14.
\begin{enumerate}
\item Decode one hexadecimal configuration string with bytes.fromhex.
\item Require sixteen bytes and return them, or raise ValueError. This is a raw-block parser, not a conversion of arbitrary text with padding.
\end{enumerate}
\subsection{\texttt{LessonTerminal.changed}}
\begin{lightcode}
def changed(self, original, faulty):
\end{lightcode}
\noindent\textbf{Source:} complete \texttt{run.py} listing in Part 10, lines 18--19.
\begin{enumerate}
\item Compare each byte of the displayed faulty value with the byte at the same position in the original.
\item Format the byte in hexadecimal, using inherited paint with red only when it differs, and return one joined display line. The byte arrays themselves remain untouched.
\end{enumerate}
\subsection{\texttt{LessonTerminal.pair}}
\begin{lightcode}
def pair(self, correct, faulty):
\end{lightcode}
\noindent\textbf{Source:} complete \texttt{run.py} listing in Part 10, lines 21--25.
\begin{enumerate}
\item Print clean bytes and a red-highlighted faulty line generated by changed.
\item XOR the two sequences and display the difference relative to an all-zero sequence. This flat view is also suitable for SM4 word-order outputs.
\end{enumerate}
\subsection{\texttt{aes\_\allowbreak{}demo}}
\begin{lightcode}
def aes_demo(terminal, limit, seed):
\end{lightcode}
\noindent\textbf{Source:} complete \texttt{run.py} listing in Part 10, lines 28--96.
\begin{enumerate}
\item Read the configured AES key and plaintext, obtain the clean round-nine trace and print normal operations and keys. For each column, inject distinct reproducible round-nine faults and intersect key\_candidates quartets.
\item Assemble all unique groups into K10, show inverse schedule steps and verify K0 by re-encryption and ground truth. Return False on ambiguity and True after successful full recovery.
\end{enumerate}
\subsection{\texttt{aes\_\allowbreak{}demo.reverse\_\allowbreak{}step}}
\begin{lightcode}
def reverse_step(index, current, transformed, recovered):
\end{lightcode}
\noindent\textbf{Source:} complete \texttt{run.py} listing in Part 10, lines 87--90.
\begin{enumerate}
\item Receive one AES schedule-reversal event and print the current word XOR transformed predecessor equals the recovered earlier word. This is an observer of reversal, not an alternate reversal algorithm.
\end{enumerate}
\subsection{\texttt{inverse\_\allowbreak{}rows}}
\begin{lightcode}
def inverse_rows():
\end{lightcode}
\noindent\textbf{Source:} complete \texttt{run.py} listing in Part 10, lines 99--101.
\begin{enumerate}
\item Accept no arguments in this compatibility adapter.
\item Return piret.inverse\_rows, preserving the shared cached inversion of the SM4 data diffusion matrix.
\end{enumerate}
\subsection{\texttt{inverse\_\allowbreak{}linear}}
\begin{lightcode}
def inverse_linear(value):
\end{lightcode}
\noindent\textbf{Source:} complete \texttt{run.py} listing in Part 10, lines 104--106.
\begin{enumerate}
\item Accept one SM4 word difference.
\item Return piret.inverse\_linear(value). The terminal module introduces no second inverse-L implementation.
\end{enumerate}
\subsection{\texttt{sm4\_\allowbreak{}encrypt}}
\begin{lightcode}
def sm4_encrypt(plaintext, key, byte=None, mask=0, fault_round=31):
\end{lightcode}
\noindent\textbf{Source:} complete \texttt{run.py} listing in Part 10, lines 109--111.
\begin{enumerate}
\item Accept an explicit plaintext/key pair and optional byte, mask and fault round.
\item Delegate to piret.sm4\_fault\_encrypt and return ciphertext plus diagnostic trace. No fault byte means a normal observed encryption.
\end{enumerate}
\subsection{\texttt{sm4\_\allowbreak{}byte\_\allowbreak{}candidates}}
\begin{lightcode}
def sm4_byte_candidates(correct, faulty, byte, mask):
\end{lightcode}
\noindent\textbf{Source:} complete \texttt{run.py} listing in Part 10, lines 114--116.
\begin{enumerate}
\item Accept a correct/faulty view and known byte/mask descriptor.
\item Return the shared piret.sm4\_byte\_candidates result. This wrapper keeps older callers compatible without duplicating candidate mathematics.
\end{enumerate}
\subsection{\texttt{reverse\_\allowbreak{}sm4}}
\begin{lightcode}
def reverse_sm4(final_keys, observer=None):
\end{lightcode}
\noindent\textbf{Source:} complete \texttt{run.py} listing in Part 10, lines 119--121.
\begin{enumerate}
\item Accept four consecutive final subkeys and an optional schedule observer.
\item Delegate to piret.reverse\_sm4 and return its recovered master key. Ordering and completeness validation belong to that shared implementation.
\end{enumerate}
\subsection{\texttt{sm4\_\allowbreak{}demo}}
\begin{lightcode}
def sm4_demo(terminal, limit, start_byte=0):
\end{lightcode}
\noindent\textbf{Source:} complete \texttt{run.py} listing in Part 10, lines 124--159.
\begin{enumerate}
\item Read the configured SM4 target, print a clean trace and define an oracle returning only faulty ciphertexts at the requested round.
\item Request full recovery using known masks one through the configured limit, display each observation and peeled view, and verify the recovered key. Return False on incomplete recovery and True on success.
\end{enumerate}
\subsection{\texttt{sm4\_\allowbreak{}demo.oracle}}
\begin{lightcode}
def oracle(number, byte, mask):
\end{lightcode}
\noindent\textbf{Source:} complete \texttt{run.py} listing in Part 10, lines 136--137.
\begin{enumerate}
\item Capture target plaintext/key, run sm4\_encrypt for the requested round, byte and mask and return only ciphertext. The engine may receive a secret key internally, but the recovery solver does not.
\end{enumerate}
\subsection{\texttt{sm4\_\allowbreak{}demo.report}}
\begin{lightcode}
def report(event):
\end{lightcode}
\noindent\textbf{Source:} complete \texttt{run.py} listing in Part 10, lines 139--146.
\begin{enumerate}
\item Receive an SM4 history dictionary and display the actual and peeled ciphertext comparisons, inverse-L difference and survivor values. It neither reads reference subkeys nor changes candidate intersections.
\end{enumerate}
\subsection{\texttt{main}}
\begin{lightcode}
def main(argv=None):
\end{lightcode}
\noindent\textbf{Source:} complete \texttt{run.py} listing in Part 10, lines 162--185.
\begin{enumerate}
\item Parse algorithm selection, stepping, colors, observation limit and seed; validate all configured raw blocks.
\item Construct the lesson terminal and execute the selected AES and/or SM4 demonstration. Return zero only when every requested lesson succeeds; the module guard also handles interruption without claiming success.
\end{enumerate}
